% =============================================================================
%  IIIT RAICHUR -- B.Tech PROJECT THESIS
%  Department of Computer Science and Engineering
%  Individual Report: Shresth Sarthak Awasthi (CS22B1051)
%  Focus: Reinforcement Learning aspects
% =============================================================================

\documentclass[12pt, a4paper, oneside]{report}

% ---------------------- PACKAGES ---------------------------------------------
\usepackage[utf8]{inputenc}
\usepackage[T1]{fontenc}
\usepackage{lmodern}
\usepackage[margin=1in, top=1.2in, bottom=1.2in]{geometry}
\usepackage{graphicx}
\usepackage{setspace}
\usepackage{titlesec}
\usepackage{tocloft}
\usepackage{fancyhdr}
\usepackage{booktabs}
\usepackage{longtable}
\usepackage{array}
\usepackage{tabularx}
\usepackage{multirow}
\usepackage{caption}
\usepackage{subcaption}
\usepackage{amsmath, amssymb}
\usepackage{xcolor}
\usepackage[hidelinks, breaklinks=true]{hyperref}
\usepackage{cleveref}
\usepackage{url}
\usepackage{enumitem}
\usepackage{float}
\usepackage[numbers]{natbib}

% ---------------------- USER INPUTS ------------------------------------------
\newcommand{\thesisTitle}{A Multi-Agent Reinforcement Learning Framework for Cooperative Non-Player Characters (NPCs) in a Role-Playing Game}
\newcommand{\thesisTitleDisplay}{A MULTI-AGENT REINFORCEMENT LEARNING FRAMEWORK \\ FOR COOPERATIVE NON-PLAYER CHARACTERS (NPCs) \\ IN A ROLE-PLAYING GAME}
\newcommand{\studentName}{Shresth Sarthak Awasthi}
\newcommand{\rollNumber}{CS22B1051}
\newcommand{\supervisorName}{Dr. Dubacharla Gyaneshwar}
\newcommand{\supervisorNameB}{Dr. Karthik P.N.}
\newcommand{\departmentName}{Computer Science and Engineering}
\newcommand{\submissionMonthYear}{May 2026}
\newcommand{\submissionLocation}{Raichur -- 584135}
\newcommand{\degreeName}{Bachelor of Technology}
\newcommand{\logoPath}{images/iiitr_logo.png}

% ---------------------- FORMATTING -------------------------------------------
\onehalfspacing

\titleformat{\chapter}[display]
  {\normalfont\huge\bfseries}{\chaptertitlename\ \thechapter}{20pt}{\Huge}
\titlespacing*{\chapter}{0pt}{-30pt}{30pt}

\pagestyle{fancy}
\fancyhf{}
\fancyfoot[C]{\thepage}
\renewcommand{\headrulewidth}{0pt}
\renewcommand{\footrulewidth}{0pt}

\fancypagestyle{plain}{%
  \fancyhf{}%
  \fancyfoot[C]{\thepage}%
  \renewcommand{\headrulewidth}{0pt}%
  \renewcommand{\footrulewidth}{0pt}%
}

\renewcommand{\cftchapfont}{\bfseries}
\renewcommand{\cftchappagefont}{\bfseries}

\newcommand{\frontmatterheading}[1]{%
  \begin{center}\Large\bfseries #1\end{center}
  \vspace{0.5cm}
}

\setlength{\parskip}{0.5em}
\setlength{\parindent}{1.5em}

% =============================================================================
\begin{document}

% ---------------------- TITLE PAGE -------------------------------------------
\begin{titlepage}
    \centering
    \thispagestyle{empty}
    \vspace*{0.2cm}

    {\Large\bfseries \thesisTitleDisplay \par}
    \vspace{0.5cm}

    {\normalsize A Project Report Submitted \par}
    {\normalsize in Partial Fulfilment of the Requirements \par}
    {\normalsize for the Degree of \par}
    \vspace{0.25cm}

    {\large\bfseries \MakeUppercase{\degreeName} \par}
    \vspace{0.3cm}

    {\normalsize in \par}
    \vspace{0.2cm}

    {\large\bfseries \MakeUppercase{\departmentName} \par}
    \vspace{0.25cm}

    {\normalsize\itshape by \par}
    \vspace{0.3cm}

    {\normalsize\bfseries \studentName{} (Roll No.~\rollNumber) \par}
    \vspace{0.25cm}

    \includegraphics[width=1\textwidth]{\logoPath}\par
    \vspace{0.25cm}

    {\normalsize\itshape to \par}
    \vspace{0.2cm}

    {\normalsize\bfseries DEPARTMENT OF \MakeUppercase{\departmentName} \par}
    {\normalsize\bfseries INDIAN INSTITUTE OF INFORMATION TECHNOLOGY \par}
    {\normalsize\bfseries RAICHUR-584135, INDIA \par}
    \vspace{0.3cm}

    {\normalsize\itshape \submissionMonthYear \par}

\end{titlepage}

% ---------------------- FRONT MATTER (roman page numbers) --------------------
\pagenumbering{roman}
\setcounter{page}{2}

% ---------------------- DECLARATION ------------------------------------------
\clearpage
\thispagestyle{plain}
\frontmatterheading{DECLARATION}
\addcontentsline{toc}{chapter}{Declaration}

\noindent I, \textbf{\MakeUppercase{\studentName}} (Roll Number: \rollNumber),
affirm that the report titled ``\textit{\thesisTitle}'', which I have submitted
to the Indian Institute of Information Technology Raichur in partial fulfilment
of the \degreeName{} programme in \departmentName{}, documents work that I
conducted myself under the mentorship of \supervisorName{} and
\supervisorNameB{}. None of the material presented here has previously earned
me, or anyone else, a degree or diploma at this or any other institution or
university. Throughout the work I have made every effort to maintain academic
integrity, and every external idea, statement, or finding that I have drawn
upon has been acknowledged with an appropriate citation.

\vspace{2.5cm}
\noindent
\begin{minipage}[t]{0.5\textwidth}
\submissionLocation \\
\submissionMonthYear
\end{minipage}
\hfill
\begin{minipage}[t]{0.4\textwidth}
\raggedleft
\studentName \\
Roll No.~\rollNumber
\end{minipage}

% ---------------------- CERTIFICATE ------------------------------------------
\clearpage
\thispagestyle{plain}
\frontmatterheading{CERTIFICATE}
\addcontentsline{toc}{chapter}{Certificate}

\noindent We certify that the project report entitled
``\textit{\thesisTitle}'', submitted by \textbf{\studentName} (Roll
No: \rollNumber) to the
Indian Institute of Information Technology Raichur in partial fulfilment
of the \degreeName{} programme in \departmentName{}, describes work carried out
by him under our supervision. To the best of our knowledge, this work has not
been submitted for the award of any degree at any other institution.

\vspace{3cm}
\noindent
\begin{minipage}[t]{0.5\textwidth}
\submissionLocation \\
\submissionMonthYear
\end{minipage}
\hfill
\begin{minipage}[t]{0.4\textwidth}
\raggedleft
(\supervisorName) \\[1em]
(\supervisorNameB) \\
Project Supervisors
\end{minipage}

% ---------------------- ACKNOWLEDGEMENTS -------------------------------------
\clearpage
\thispagestyle{plain}
\frontmatterheading{ACKNOWLEDGEMENTS}
\addcontentsline{toc}{chapter}{Acknowledgements}

\noindent I owe my deepest thanks to my supervisors,
\textbf{\supervisorName} and \textbf{\supervisorNameB}. Their guidance, their
thoughtful critiques, and their steady encouragement carried this project
forward, and it was their mentorship that gave the work both its technical
direction and its academic discipline.

\noindent I am grateful to the Department of \departmentName{} at the Indian
Institute of Information Technology, Raichur, whose academic setting and
facilities made this work possible, and equally to the Department of Artificial
Intelligence at the Indian Institute of Technology Hyderabad, which provided
the computing resources behind the experiments.

\noindent My thanks also go to the evaluating faculty for the time they devoted
to assessing this work, and to my partner on this two-person project for a
genuinely productive collaboration. Last of all, I remain deeply indebted to my
family and friends, whose support and encouragement never wavered.

\vspace{2cm}
\noindent
\begin{flushright}
\studentName \\
Roll No.~\rollNumber
\end{flushright}

% ---------------------- ABSTRACT ---------------------------------------------
\clearpage
\thispagestyle{plain}
\frontmatterheading{ABSTRACT}
\addcontentsline{toc}{chapter}{Abstract}

This report develops a multi-agent reinforcement learning framework in which cooperative behaviour emerges among NPC agents inside a medieval-village RPG used as a research playground. Its central contribution is a decomposed reward model, built around a dynamic $\lambda$ coefficient, that couples each agent's individual tendency to cooperate with the weight it places on community welfare: $R_i(t) = P_i(t) + G_i(t) + \lambda_i(t) \cdot C(t)$, where $\lambda_i = 0.05 + \text{coop}_i \times 0.55$. Six tabular Q-learning agents act within a discretized 225-state space, each subject to role-specific action masking and adaptive personality dynamics. To probe how robustly they learn under distribution shift, a system shock engine introduces non-stationary perturbations such as famine, plague, and raids. Across an ablation study spanning six conditions of 500 turns apiece, we show that: (1)~Q-learning drives cooperation from a 0.000 baseline to 0.699 while schedule-only agents show no emergence; (2)~role-specific action masking improves cooperation by 0.161 through structural inductive bias, the largest single contributor after RL itself; (3)~the famine shock has a substantial, lasting effect --- removing shocks raises final cooperation by 0.237 (to 0.936) --- yet agents keep learning rather than collapsing; (4)~the dynamic $\lambda$ feedback loop amplifies cooperation by a further 0.104 over a static baseline; and (5)~adaptive personality coefficients produce meaningful role divergence in social sensitivity and risk aversion. Beyond the learning core, the system offers optional LLM narration backed by template fallbacks, and its game world is organised as a hierarchical MDP that drives quest progression.

% ---------------------- TABLE OF CONTENTS ------------------------------------
\clearpage
\tableofcontents
\listoffigures
\addcontentsline{toc}{chapter}{List of Figures}
\listoftables
\addcontentsline{toc}{chapter}{List of Tables}

\clearpage
\pagenumbering{arabic}
\setcounter{page}{1}

% =============================================================================
% CHAPTER 1: INTRODUCTION
% =============================================================================
\chapter{Introduction}
\label{chap:introduction}

In most modern games, non-player characters (NPCs) run on scripted behaviour trees or finite state machines, which makes their interactions predictable and repetitive. This project asks whether reinforcement learning can instead give rise to emergent cooperative behaviour in a multi-agent game, without cooperation ever being scripted explicitly.

We build a single-player medieval village RPG where six NPCs learn autonomously via tabular Q-learning. The central research question is:

\begin{quote}
\textit{Can NPCs learn to shift from individually optimal strategies to cooperative behavior when community welfare contributes to their total reward?}
\end{quote}

\section{Problem Statement}
\label{sec:problem-statement}

Today's in-game NPC systems depend on hand-authored behaviour trees, finite state machines, or scripted dialogue. Such approaches yield deterministic, repetitive behaviour that breaks immersion and fails to adapt to how the player plays. The central problems are:

\begin{itemize}
    \item \textbf{No emergent cooperation:} Scripted NPCs cannot discover cooperative strategies on their own --- any prosocial behavior must be explicitly programmed by designers for every scenario.
    \item \textbf{Brittleness to perturbation:} Static behavior trees cannot adapt when the game environment changes (e.g., resource shocks, new player strategies). NPCs continue executing their scripts regardless of context.
    \item \textbf{No individual--community trade-off:} Real social agents must balance personal gain against group welfare. Scripted NPCs have no mechanism to represent, let alone dynamically adjust, this trade-off.
    \item \textbf{Lack of interpretability in deep RL alternatives:} While deep RL can produce adaptive NPC behavior, the learned policies are opaque. For a research playground, full interpretability of what agents have learned is essential.
\end{itemize}

\section{Project Significance}
\label{sec:project-significance}

This project contributes to both game AI and multi-agent reinforcement learning research:

\begin{itemize}
    \item \textbf{Emergent social dynamics:} Demonstrates that cooperation can arise purely from reward structure and learning, without being scripted --- a result relevant to computational social science and mechanism design.
    \item \textbf{Interpretable multi-agent RL:} By using tabular Q-learning (6,300 Q-values per NPC), every learned policy is directly inspectable. This contrasts with deep RL approaches where policy interpretation requires post-hoc explanation methods.
    \item \textbf{Decomposed reward framework:} The $R_i(t) = P_i(t) + G_i(t) + \lambda_i(t) \cdot C(t)$ formulation with dynamic $\lambda$ is domain-agnostic and transferable to industrial multi-agent coordination problems (factory automation, automotive ECU tuning).
    \item \textbf{Research playground:} The system provides a complete experimental platform with ablation controls, analytics dashboard, deterministic seeding, and automated bot playthroughs for reproducible research.
\end{itemize}

\section{Scope of Work}
\label{sec:scope}

The system is structured in three layers:

\begin{itemize}
    \item \textbf{Hierarchical MDP:} A 7-stage quest graph with 16 static checkpoints drives game progression. Dynamic checkpoints handle player deviations with nudging and convergence mechanisms.
    \item \textbf{Q-Learning NPC Agents:} Six NPCs run tabular Q-learning over a 225-state $\times$ 28-action space with role-specific action masking and adaptive personality coefficients.
    \item \textbf{Optional LLM Narration:} An HTTP adapter connects to OpenRouter, Ollama, or any OpenAI-compatible provider. Every LLM call has a deterministic template fallback, ensuring the game runs identically without a model.
\end{itemize}

The project is implemented as a web application (Python/FastAPI backend, vanilla JavaScript frontend) with a real-time analytics dashboard for monitoring cooperation curves, reward decomposition, and shock response.

\section{Objectives}
\label{sec:objectives}

The main objectives of this project are:

\begin{enumerate}
    \item Design and implement a multi-agent reinforcement learning environment within a medieval village RPG setting.
    \item Develop a decomposed reward model with a dynamic $\lambda$ coefficient that creates feedback loops between individual and community welfare.
    \item Implement six tabular Q-learning NPC agents with role-specific action masking and adaptive personality dynamics.
    \item Build a system shock engine to test agent resilience under non-stationary environmental perturbations.
    \item Conduct a systematic ablation study isolating the contribution of each subsystem to emergent cooperation.
    \item Integrate optional LLM narration with guardrails and template fallbacks for atmospheric game text.
\end{enumerate}

\section{Novelty}
\label{sec:novelty}

The key novel contributions of this project are:

\begin{itemize}
    \item \textbf{Dynamic $\lambda$ feedback loop:} Unlike static reward weighting in prior MARL work, the $\lambda_i(t) = 0.05 + \text{coop}_i \times 0.55$ mechanism creates a virtuous cycle where cooperation self-amplifies through the reward structure. Our ablation shows this contributes $+0.104$ to cooperation --- one of the largest factors, alongside role masking, after RL itself.
    \item \textbf{Fully interpretable multi-agent cooperation:} While prior game AI research uses deep RL (opaque policies), our tabular approach lets researchers directly inspect every Q-value and trace exactly why an NPC chose a cooperative action.
    \item \textbf{Integrated research playground:} Within one reproducible platform, the system unifies RL agents, environmental shocks, adaptive personalities, quest progression, LLM narration, and a full analytics dashboard, all under explicit ablation controls --- an assembly we have not seen brought together in existing game-AI testbeds.
\end{itemize}

% =============================================================================
% CHAPTER 2: LITERATURE SURVEY
% =============================================================================
\chapter{Literature Survey}
\label{chap:literature}

Four distinct research areas intersect within this project, and this chapter surveys the existing literature in each: reinforcement learning across multiple agents, the shaping of reward signals, intelligent behaviour for non-player characters, and narration augmented by language models. Within every area we outline the dominant techniques and pinpoint the shortcomings that motivate our own contributions.

\section{Multi-Agent Reinforcement Learning}
\label{sec:marl}

A large part of the theory underpinning reinforcement learning goes back to Watkins and Dayan~\cite{watkins1992qlearning}. They showed that, when every state-action pair is revisited without bound and the learning rate decays appropriately, tabular Q-learning converges on an optimal policy inside a single-agent MDP. For a broad treatment of the value-based methods, policy-gradient formulations, and temporal-difference algorithms on which modern RL rests, the reader may consult Sutton and Barto~\cite{sutton2018rl}.

Once several agents share an environment, the problem grows considerably harder. The survey of Busoniu et al.~\cite{busoniu2008marl} identifies three core difficulties. The first is non-stationarity: from any one agent's vantage point the environment keeps shifting, because every other agent is busy updating its own policy at the same time. The second is the combinatorial blow-up of the joint action space, which grows exponentially with the number of agents. The third is credit assignment --- when a single team reward is shared, it is far from obvious which agent's action actually earned it. On the convergence side, Bowling and Veloso~\cite{bowling2002multiagent} propose adjusting the learning rate depending on whether an agent is winning or losing; the resulting guarantees, however, cover only two-player games and do not extend cleanly to the general $N$-agent setting.

The simplest of the multi-agent approaches is the independent-learner (IL) scheme, in which each agent runs an ordinary single-agent algorithm and treats all the others as part of its environment. This violates the stationarity assumption outright, yet IL tends to work surprisingly well in practice as long as the agents change slowly relative to their own learning rate~\cite{busoniu2008marl}. Our design adopts exactly this IL philosophy: every NPC owns a separate tabular Q-learning instance. Since adaptation is deliberately slow (rate 0.018) and a homeostatic pull nudges the coefficients back toward 0.5, the two act as built-in regularisers and keep the environment looking approximately stationary to each individual agent.

Richer frameworks have since been built for cooperative MARL. QMIX~\cite{rashid2018qmix} factors the joint Q-value under a monotonicity constraint, which supports centralized training together with decentralized execution on cooperative tasks. MADDPG~\cite{lowe2017maddpg} extends actor-critic methods to the multi-agent case by letting each agent's critic see the observations of all agents during training. In grid-world settings, Leibo et al.~\cite{leibo2017multi} study how social dilemmas arise and show that whether independent learners end up cooperating or defecting depends on the incentive structure of the environment.

By coupling independent Q-learners to a shared community reward, our approach lies closest to the emergent-cooperation line of Leibo et al., but with one decisive addition: a dynamic $\lambda$ coefficient that tunes how strongly each agent listens to the community signal based on its own past record of cooperation. This gives the agents an implicit means of coordinating, all while doing away with the centralized critics that QMIX and MADDPG require.

\section{Reward Shaping}
\label{sec:reward-shaping}

Ng et al.~\cite{ng1999rewardshaping} proved that potential-based reward shaping accelerates learning while leaving the set of optimal policies untouched. Their landmark result states that a shaping term written as $F(s, s') = \gamma \Phi(s') - \Phi(s)$ changes only how quickly the agent learns, never the policy it finally arrives at. The key caveat is that, if the shaping reward is not cast as a difference of potentials, it can create spurious equilibria.

Our decomposed reward formulation puts this principle to work in a specific way: the community term $\lambda_i(t) \cdot C(t)$ acts as a shaping signal that pushes agents toward prosocial behaviour, while the dynamic $\lambda$ sets the strength of that push according to each agent's cooperation history. Strictly, the scheme is not potential-based --- because $\lambda$ varies with the agent's state, the shaping function becomes non-stationary --- yet its behaviour in practice is close to that of potential-based shaping. An agent that begins to cooperate is pulled toward even greater cooperation by the $\lambda$ feedback loop, which speeds convergence to cooperative equilibria. The ablation makes this acceleration tangible: C7 with a static $\lambda$ settles at 0.595, whereas C1 with the dynamic $\lambda$ climbs to 0.699 --- a direct measure of what adaptive shaping buys over a fixed-shaping baseline.

\section{NPC Intelligence in Games}
\label{sec:npc-intelligence}

In commercial titles, NPC behaviour is conventionally controlled by behaviour trees~\cite{isla2005bt}, finite state machines, or goal-oriented action planning --- techniques that force designers to author every decision branch by hand. Despite decades of RL research, its adoption in shipped commercial games remains thin~\cite{liu2025schola}. After interviewing 17 creators of game agents across both AAA and indie studios, Jacob et al.~\cite{jacob2020unwieldy} found that the industry avoids RL because the learned behaviour is hard to predict, training eats up time, and the tooling around it is lacking. Yannakakis and Togelius~\cite{yannakakis2018gameai} make a related point, noting that commercial games have largely passed over learning-based AI. On the rare occasions it does surface, it stays within narrow, tightly scoped subtasks --- placing a goalkeeper, say, or capturing a driving style --- rather than steering general NPC decisions or social conduct.

Although recent academic studies have investigated RL for game NPCs~\cite{riedl2012interactive}, these efforts have largely targeted single-agent settings. As far as we are aware, no earlier study brings \textit{cooperative} multi-agent RL together with dynamic reward weighting in a complete game world that combines quest progression, social interaction among several NPCs, and environmental disturbances.

\section{LLM-Enhanced Game Narration}
\label{sec:llm-narration}

Interactive fiction~\cite{urbanek2019light} and the generation of in-game dialogue are among the tasks to which large language models have been put. In our system, by contrast, an LLM serves merely as an optional narration layer governed by a nine-stage guardrail pipeline; because game logic is kept distinct from narration, the game remains fully playable even when no model is available.

\section{Comparison with Our Approach}
\label{sec:lit-comparison}

\begin{table}[H]
\centering
\caption{Comparison of Related Approaches and Our System}
\label{tab:lit-comparison}
\small
\begin{tabularx}{\textwidth}{|l|l|X|X|}
\hline
\textbf{Work} & \textbf{Year} & \textbf{Methodology} & \textbf{Our Differentiation} \\
\hline
QMIX~\cite{rashid2018qmix} & 2018 & Joint Q-value decomposition under central training & Independent learners sharing a community reward, removing any need for a central critic \\
\hline
MADDPG~\cite{lowe2017maddpg} & 2017 & Central critic with decentralized actors (deep RL) & Fully interpretable tabular Q-learning; dynamic $\lambda$ replaces static reward mixing \\
\hline
Leibo et al.~\cite{leibo2017multi} & 2017 & Emergent social dilemmas studied in grid-worlds & Complete RPG world featuring quests, shocks, adaptive personalities, and LLM narration \\
\hline
Ng et al.~\cite{ng1999rewardshaping} & 1999 & Theory of potential-based reward shaping & Augmented with dynamic gating through $\lambda_i(t)$ that tracks cooperation tendency \\
\hline
Watkins \& Dayan~\cite{watkins1992qlearning} & 1992 & Convergence proof for tabular Q-learning & Carried into a multi-agent game with role-specific masking and adaptive personalities \\
\hline
Bowling \& Veloso~\cite{bowling2002multiagent} & 2002 & Adaptive learning rate for multi-agent convergence & Slow-adapting independent learners as implicit regularisation against non-stationarity \\
\hline
Isla~\cite{isla2005bt} & 2005 & Behaviour trees governing NPC AI & Hand-authored trees give way to RL; cooperation arises instead of being scripted \\
\hline
LIGHT~\cite{urbanek2019light} & 2019 & Interactive fiction driven by an LLM & LLM acts only as optional narration, not the engine; nine-stage guardrails \\
\hline
\end{tabularx}
\end{table}

% =============================================================================
% CHAPTER 3: METHODOLOGY
% =============================================================================
\chapter{Methodology}
\label{chap:methodology}

This chapter details the technical design of the system, with emphasis on the reinforcement learning architecture: the Q-learning formulation, decomposed reward model with dynamic $\lambda$, action masking, adaptive personality dynamics, and the system shock engine that tests learning robustness. The hierarchical MDP for quest progression and the LLM integration layer are covered more briefly.

\section{Hierarchical MDP}
\label{sec:hierarchical-mdp}

The quest is modeled as a two-level hierarchical MDP. At the macro level, seven stages drive game progression:

\begin{equation}
    S_1 \to S_2 \to S_3 \to S_4 \to S_5 \to S_6 \to S_7 \to S_{\text{success}}
\end{equation}

Every macro stage $S_k$ encloses a micro MDP $\mathcal{M}_{\text{micro}}^{(k)}$ carrying its own checkpoints, completion conditions, and transition rewards. The seven stages cover arrival at the gate, the search for the elder, the unveiling of the quest, investigation, retrieval, and the return journey.

When the player strays from the intended path, the system spawns dynamic checkpoints that register the detour and steer the player back toward the nearest static checkpoint. If the quest is finished before the turn budget runs out, its state is reset while every NPC's Q-table and adaptation state is retained, which lets learning continue across multiple episodes inside one session. This auto-restart behaviour is essential to the RL results, because it lets the NPCs build up learning experience over several quest cycles within a single 500-turn run.

\section{NPC State Space and Q-Learning}
\label{sec:q-learning}

Each NPC maintains a tabular Q-table over a discretized state space:

\begin{table}[H]
\centering
\caption{NPC State Space Dimensions}
\label{tab:state_space}
\begin{tabular}{lll}
\toprule
\textbf{Dimension} & \textbf{Values} & \textbf{Count} \\
\midrule
Location & gate, village\_center, elders\_house, fields, tavern & 5 \\
Time & morning, midday, afternoon, evening, night & 5 \\
Energy & low ($<$33\% HP), medium (33--66\%), high ($>$66\%) & 3 \\
Mood & unhappy ($<$4), content (4--7), cheerful ($>$7) & 3 \\
\bottomrule
\end{tabular}
\end{table}

This yields a total state space of $5 \times 5 \times 3 \times 3 = 225$ states. The action space holds 28 universal actions grouped into 7 categories: navigation, exploration, talk, social, combat, stealth, and utility. Each NPC therefore stores a Q-table of 6,300 entries.

\textbf{Q-Learning update:}
\begin{equation}
    Q(s, a) \leftarrow Q(s, a) + \alpha \left[ r + \gamma \max_{a'} Q(s', a') - Q(s, a) \right]
\end{equation}
Here $\alpha = 0.1$ and $\gamma = 0.9$, and $\varepsilon$-greedy exploration decays from 0.15 to 0.05 at a per-turn rate of $0.995$. Each NPC is pre-trained for 100 episodes of 50 turns at $\varepsilon = 0.5$.

\textbf{Why tabular Q-learning over deep RL.}
Choosing tabular Q-learning instead of a function approximator such as DQN was a deliberate decision resting on three considerations. The first is \emph{interpretability}: because each NPC holds only 6,300 Q-values, the whole learned policy is open to inspection --- one can read off the Q-value of any state-action pair and follow precisely why a given action was favoured. Deep RL policies, on the other hand, need post-hoc tools such as saliency maps or SHAP values, which approximate rather than exactly recover that attribution. The second is \emph{convergence}: Watkins and Dayan~\cite{watkins1992qlearning} established that tabular Q-learning reaches the optimal policy under standard conditions. The multi-agent setting does break the stationarity assumption, but the slow adaptation rate together with the homeostatic drift yields empirical stability, borne out by the steady convergence to $\mathcal{K} \approx 0.699$ across runs. The third is \emph{tractability}: a 225-state space is small enough that, over a 500-turn experiment with quest auto-restart, every state-action pair is visited repeatedly, so the Q-values are properly estimated rather than interpolated.

\textbf{Epsilon-greedy exploration strategy.}
Action selection follows an $\varepsilon$-greedy rule whose exploration rate decays as $\varepsilon_t = \max(0.05,\; 0.15 \times 0.995^t)$, keeping exploration and exploitation in balance over the course of the experiment. The starting value of $\varepsilon = 0.15$ leaves enough room to uncover cooperative strategies during the early turns, and the exponential decay brings the rate down to roughly 0.055 by the midpoint ($t \approx 200$), at which stage the agents largely exploit what they have learned. A floor of 0.05 preserves a slim chance of a random move throughout, which guards against settling too early into a poor local equilibrium. Pre-training runs at a far higher $\varepsilon = 0.5$ so that the state-action space is broadly covered, equipping each NPC with a sensible initial Q-table before the main experiment starts.

\section{Decomposed Reward Model}
\label{sec:reward-model}

Each NPC $i$ receives a reward signal composed of three terms:
\begin{equation}
    R_i(t) = P_i(t) + G_i(t) + \lambda_i(t) \cdot C(t)
    \label{eq:reward}
\end{equation}

\begin{itemize}[nosep]
    \item $P_i(t)$: \textbf{Penalty} --- $-10$ when HP drops below 20\%.
    \item $G_i(t)$: \textbf{Individual reward} --- $\sum_k w_k \cdot (\text{stat}_k^{\text{new}} - \text{stat}_k^{\text{old}})$, weighted by archetype-specific preferences (e.g., farmers weight income at 0.4, guards weight reputation at 0.4).
    \item $C(t)$: \textbf{Community reward} --- Village-level welfare computed as $C(t) = 0.5 \cdot C_{\text{abs}}(t) + 0.5 \cdot C_{\Delta}(t)$, using centered non-linear shaping (sigmoid for reputation and mood, sub-linear power for health) that can go negative when village welfare drops.
    \item $\lambda_i(t)$: \textbf{Dynamic lambda} --- $\lambda_i = 0.05 + \text{coop}_i \times 0.55$, ranging from 0.05 (selfish) to 0.60 (cooperative). Creates a feedback loop: positive community reward $\to$ cooperation rises $\to$ $\lambda$ increases $\to$ more weight on community reward.
\end{itemize}

\textbf{Intuition behind the three components.}
The decomposition into $P$, $G$, and $\lambda \cdot C$ is designed to capture three distinct pressures that real social agents face. The penalty term $P_i(t)$ acts as a hard survival constraint: no matter how cooperative an agent becomes, it must avoid actions that endanger its own existence. This prevents degenerate cooperation where agents sacrifice themselves for the group. The individual reward $G_i(t)$ captures role-specific self-interest: a farmer that never improves its income is failing at its primary function, regardless of village welfare. The community term $\lambda_i(t) \cdot C(t)$ introduces the social dimension, but critically, it is \emph{gated} by the agent's own cooperation history via $\lambda$. This means the community signal is not imposed uniformly; agents that have historically cooperated receive a stronger community gradient, while agents that have not cooperated remain primarily self-interested. The interaction between these three terms creates a natural tension: $G_i$ pulls agents toward individually optimal actions, $\lambda_i \cdot C$ pulls toward prosocial actions, and $P_i$ provides a floor below which no agent will go. The dynamic $\lambda$ resolves this tension over time by amplifying whichever direction the agent has been moving.

\section{Role-Specific Action Masking}
\label{sec:action-masking}

Each NPC archetype has a preferred action set. During action selection:
\begin{itemize}[nosep]
    \item \textbf{Hard masking}: Invalid actions get $Q(s,a) = -\infty$.
    \item \textbf{Soft masking}: Role-aligned actions receive a $+0.5$ Q-value bonus; misaligned actions receive a $-0.3$ penalty.
\end{itemize}

Soft masking functions as an inductive bias that steers exploration without shutting it down. Take a guard NPC posted at the gate: its combat and patrol actions pick up the $+0.5$ bonus, so $\varepsilon$-greedy exploration is more likely to select them. Even so, the agent can still find that a social action pays off in community reward for particular states --- the $-0.3$ penalty is small next to the Q-value range and a strong reward signal can override it. The ablation outcome (C6 vs.\ C1, $\Delta = -0.161$) bears out that masking speeds up learning by shrinking the effective exploration space, giving cleaner Q-value estimates with less noise from irrelevant actions --- in fact, masking is the single largest contributor to cooperation after RL itself.

\section{Adaptive Personality Dynamics}
\label{sec:adaptation}

Four adaptation coefficients evolve per turn at base rate 0.018 with homeostatic drift rate 0.004 toward 0.5:

\begin{table}[H]
\centering
\caption{Adaptation Coefficients}
\label{tab:adaptation}
\begin{tabular}{lll}
\toprule
\textbf{Coefficient} & \textbf{Drives} & \textbf{Role Multipliers} \\
\midrule
cooperation\_tendency & $\lambda$ (community weight) & Gated by $\lambda$: $r_{\text{coop}} = r_{\text{base}} \times (0.5 + \lambda_i)$ \\
risk\_aversion & Risk avoidance & Guard 1.5$\times$, Farmer 1.2$\times$, Elder 0.6$\times$ \\
social\_sensitivity & Social feedback & Elder 1.4$\times$, Tavern 1.2$\times$, Guard 0.5$\times$ \\
shock\_resilience & Shock adaptation & Increases under shock pressure \\
\bottomrule
\end{tabular}
\end{table}

\textbf{Homeostatic drift as a stability mechanism.}
The homeostatic drift rate of 0.004 toward 0.5 serves a critical role in preventing runaway coefficient dynamics. Without it, the $\lambda$ feedback loop could drive cooperation\_tendency to its maximum (1.0) and lock agents into inflexible cooperative behavior, or conversely, a sequence of negative shocks could drive it to its minimum (0.0) and permanently suppress cooperation. The drift acts as a slow mean-reversion force: at each turn, each coefficient is pulled 0.004 toward the neutral midpoint of 0.5. This means that sustained positive or negative pressure is required to maintain extreme coefficient values. If the environmental pressure disappears (e.g., after a shock expires), the coefficients gradually return to baseline, restoring the agent's capacity to respond to future perturbations. The interaction between the adaptation rate (0.018) and the drift rate (0.004) creates a net rate of approximately 0.014 when the agent receives consistent signals, slow enough to avoid oscillation but fast enough to produce meaningful divergence over 500 turns.

\section{System Shock Engine}
\label{sec:shock-engine}

Five shock types inject non-stationary perturbations:

\begin{table}[H]
\centering
\caption{System Shock Catalog}
\label{tab:shocks}
\begin{tabular}{llll}
\toprule
\textbf{Shock} & \textbf{Duration} & \textbf{Reward Scale} & \textbf{Purpose} \\
\midrule
Famine & 15 turns & 0.5$\times$ & Resource scarcity \\
Bandit Raid & 8 turns & 0.6$\times$ & External threat \\
Plague & 12 turns & 0.4$\times$ & Sustained decline \\
Trade Boom & 10 turns & 1.5$\times$ & Recovery (positive) \\
Harsh Winter & 20 turns & 0.7$\times$ & Long-duration stress \\
\bottomrule
\end{tabular}
\end{table}

Shocks exert direct downward pressure on cooperation: once a shock's intensity passes 0.1, cooperation tendency falls in proportion to it, regardless of what the reward signal is doing.

\textbf{Non-stationarity and Q-learning robustness.}
For Q-learning convergence, system shocks pose a particular problem. The convergence argument of Watkins and Dayan~\cite{watkins1992qlearning} presumes a stationary MDP, yet a shock briefly rewrites both the reward function (through the reward-scale factor) and the transition dynamics (through stat drains and trust modifiers). As a result, Q-values acquired before the shock turn partly stale while it lasts. The tabular update, however, is inherently local --- it touches only the $Q(s,a)$ entry of the visited state-action pair --- and with the learning rate at $\alpha = 0.1$ those stale values are corrected bit by bit. The empirical comparison (C5 vs.\ C1, $\Delta = +0.237$) reveals that the famine shock leaves a sizeable and durable mark on learning: it carves a clear dip in the community-reward signal at turn 166 (\cref{fig:reward}) and holds cooperation down relative to the shock-free run (C5 attains 0.936 against C1's 0.699). What matters is that the agents do not collapse --- the local update, aided by the shock\_resilience coefficient that climbs under pressure, keeps cooperation rising once the shock lifts --- but within the 500-turn session they never claw back the lost ground, so one famine leaves a lasting, compounding imprint by way of the $\lambda$ feedback loop.

\section{LLM Integration}
\label{sec:llm-integration}

The LLM service takes the form of a lightweight HTTP adapter that works with Ollama, OpenRouter, and any OpenAI-compatible provider. It covers four use cases --- input parsing, NPC dialogue, checkpoint generation, and narration enhancement --- and its output is policed by a nine-stage guardrail pipeline that strips reasoning blocks, parses JSON, validates the schema and enums, clamps values, checks action references, filters content, and enforces length. Because every code path falls back to a deterministic template, the RL experiments behave identically whether or not an LLM is present. The default model is \texttt{qwen3:4b} through LM Studio, or \texttt{meta-llama/llama-3.2-3b-instruct:free} through OpenRouter.

\section{Turn Processing Pipeline}
\label{sec:turn-pipeline}

Each player action triggers a 10-step pipeline:
\begin{enumerate}[nosep]
    \item Resolve player action (combat, social, navigation, etc.)
    \item Log event with importance scoring
    \item Check quest completion or deviation
    \item Process 6 NPC turns (Q-learning action selection, reward computation, Q-table update)
    \item Check random events
    \item Shock engine tick (decay intensities, apply stat drain, trust modifiers)
    \item Time advancement (4 turns = 1 period)
    \item Reputation decay, stamina regeneration, perception checks
    \item Game-over conditions (or quest restart if turns remain)
    \item Telemetry logging (world snapshot, reward decomposition, adaptation state)
\end{enumerate}

% =============================================================================
% CHAPTER 4: CHALLENGES AND SOLUTIONS
% =============================================================================
\chapter{Challenges and Solutions}
\label{chap:challenges}

Developing a multi-agent RL system within a full game environment introduced several technical and design challenges. Table~\ref{tab:challenges} summarizes the major issues faced and their corresponding solutions.

\begin{table}[H]
\centering
\caption{Major Implementation Challenges and Their Solutions}
\label{tab:challenges}
\begin{tabularx}{\textwidth}{|X|X|}
\hline
\textbf{Challenge} & \textbf{Solution} \\
\hline
\textbf{Multi-agent non-stationarity:} Six Q-learners in a shared environment violate the stationary MDP assumption needed for convergence guarantees.
& Slow adaptation rate (0.018), homeostatic drift toward 0.5, and a shared community reward signal that aligns agents, empirically producing stable cooperation. \\
\hline
\textbf{Reward hacking risk:} Agents could exploit $C(t)$ via degenerate strategies that inflate village stats without genuine cooperation.
& Competing objectives (role-specific $G_i$, penalty $P_i$) and non-linear shaping (sigmoid, sub-linear) prevent saturation exploits. Ablation verified cooperation correlates with prosocial actions. \\
\hline
\textbf{LLM output quality:} Local 4B models produce inconsistent JSON, hallucinate game state, and occasionally generate inappropriate content.
& 9-stage guardrail pipeline (JSON parse, schema/enum validation, value clamping, content filtering) with deterministic template fallbacks on every code path. \\
\hline
\textbf{Reproducibility:} Python's hash randomization and platform-specific RNG behavior make cross-run reproducibility difficult.
& Single \texttt{MASTER\_SEED}; NPC pre-training seeds derived from registry index (not \texttt{hash()}). Explicit reproducibility tests verify identical playthroughs. \\
\hline
\textbf{Windows asyncio deadlocks:} Module-based execution caused \texttt{ProactorEventLoop} deadlocks with \texttt{asyncio.to\_thread}.
& Direct script execution with \texttt{sys.path} insertion, avoiding the Windows-specific event loop interaction. \\
\hline
\end{tabularx}
\end{table}

% =============================================================================
% CHAPTER 5: RESULTS AND DISCUSSION
% =============================================================================
\chapter{Results and Discussion}
\label{chap:results}

This chapter presents the experimental design, ablation study results, reward analysis, NPC adaptation patterns, and narrative quality evaluation.

\section{Experimental Design}
\label{sec:experimental-design}

\subsection{Ablation Conditions}

We define seven ablation conditions, each isolating one subsystem:

\begin{table}[H]
\centering
\caption{Ablation Conditions}
\label{tab:conditions}
\begin{tabular}{llccccc}
\toprule
\textbf{ID} & \textbf{Description} & \textbf{MDP} & \textbf{RL} & \textbf{Shocks} & \textbf{Mask} & \textbf{Dyn.\ $\lambda$} \\
\midrule
C1 & Full System (baseline) & \checkmark & \checkmark & \checkmark & \checkmark & \checkmark \\
C2 & No LLM & \checkmark & \checkmark & \checkmark & \checkmark & \checkmark \\
C3 & No RL (schedules only) & \checkmark & --- & \checkmark & --- & --- \\
C4 & Flat MDP (no hierarchy) & --- & \checkmark & \checkmark & \checkmark & \checkmark \\
C5 & No Shocks & \checkmark & \checkmark & --- & \checkmark & \checkmark \\
C6 & No Role Masking & \checkmark & \checkmark & \checkmark & --- & \checkmark \\
C7 & Static $\lambda = 0.325$ & \checkmark & \checkmark & \checkmark & \checkmark & --- \\
\bottomrule
\end{tabular}
\end{table}

\subsection{Metrics}

\begin{itemize}[nosep]
    \item \textbf{Cooperation index} $\mathcal{K}(t) = \frac{1}{N} \sum_{i=1}^{N} \text{coop}_i(t)$ --- the mean cooperation tendency taken over all NPCs at a given turn.
    \item \textbf{Social welfare index} --- average cooperation tendency weighted by community reward performance.
    \item \textbf{Policy entropy} $\mathcal{D}_i = -\sum_{a \in \mathcal{A}} p_i(a) \log p_i(a)$ --- derived from the complete per-action counts.
    \item \textbf{Contextual grounding} --- fraction of game-state keywords (action, location, NPC names) appearing in each narration.
    \item \textbf{Vocabulary diversity} --- type-token ratio (unique words / total words) across narrations.
    \item \textbf{Shock recovery time} --- turns after shock expiry until $\mathcal{K}(t)$ returns to pre-shock baseline.
\end{itemize}

\subsection{Protocol}

All runs are automated bot playthroughs driven by the \texttt{quest\_focused} strategy under seed 42 to keep them reproducible. Every condition lasts 500 turns with quest auto-restart enabled. In conditions that include shocks, a famine is triggered at turn 166. The difficulty is fixed to the normal preset, with adaptive difficulty switched off.

\section{Ablation Study Results}
\label{sec:ablation-results}

\begin{table}[H]
\centering
\caption{Ablation Study Results (500 turns, seed=42)}
\label{tab:results}
\begin{tabular}{llcc}
\toprule
\textbf{Condition} & \textbf{Description} & \textbf{Final $\mathcal{K}$} & \textbf{$\Delta$ vs C1} \\
\midrule
C1 & Full System (baseline) & 0.699 & --- \\
C3 & No RL (schedules only) & 0.000 & $-0.699$ \\
C4 & Flat MDP (no hierarchy) & 0.871 & $+0.172$ \\
C5 & No Shocks & 0.936 & $+0.237$ \\
C6 & No Role Masking & 0.538 & $-0.161$ \\
C7 & Static $\lambda = 0.325$ & 0.595 & $-0.104$ \\
\bottomrule
\end{tabular}
\end{table}

\begin{figure}[H]
    \centering
    \includegraphics[width=\textwidth]{images/fig_cooperation.png}
    \caption{Cooperation index $\mathcal{K}(t)$ over 500 turns for six ablation conditions. C3 (No RL) stays flat at 0.000 while all RL-enabled conditions rise from zero. C5 (No Shocks) reaches the highest final value and C4 (Flat MDP) is second, both cycling through more learning episodes than C1. C6 (No Role Masking) and C7 (Static $\lambda$) rise more slowly, confirming the contributions of role masking and the dynamic $\lambda$ feedback loop. The famine shock at turn 166 lastingly depresses C1 relative to the shock-free C5.}
    \label{fig:cooperation}
\end{figure}

\subsection{Finding 1: RL Drives Cooperation}

The headline result is this: C1 (Full System) reaches a cooperation index of 0.699, whereas C3 (No RL) stays at exactly 0.000 --- the initialization baseline, since cooperation tendency begins at zero. Stripped of Q-learning, the NPCs simply execute fixed schedules and cannot adjust their behaviour. The 0.699 gap shows that cooperation \textit{emerges} from the learning process itself rather than falling out of the game mechanics --- RL is, by a wide margin, the dominant driver.

\subsection{Finding 2: Role Masking Aids Learning}

C6 (No Role Masking) lands at 0.538 against C1's 0.699 --- a gap of 0.161, which is the single largest component contribution to cooperation once RL itself is set aside. Lacking role constraints, the NPCs sweep the full 28-action space evenly, and their Q-value estimates come out noisier as a result. Role masking supplies a structural inductive bias: guards learn chiefly from guard-relevant actions and farmers from farming actions, so the gradient signals they produce are cleaner.

\subsection{Finding 3: Shock Resilience}

C5 (No Shocks) reaches 0.936, the highest of all conditions and 0.237 above C1's 0.699. The famine shock injected at turn 166 produces a visible dip in C1's community-reward signal (\cref{fig:reward}) and a lasting downward divergence in its cooperation trajectory: before the shock, C1 and C5 track together, but afterward C5 pulls steadily ahead (\cref{fig:cooperation}). The agents remain robust in that cooperation keeps rising rather than collapsing --- the shock\_resilience coefficient increases under pressure and the local Q-updates correct stale values --- but they do not recover the lost ground within the session. Shocks are thus the largest single environmental drag on cooperation in our study.

\subsection{Finding 4: MDP Structure}

C4 (Flat MDP) climbs to 0.871, the second-highest figure (trailing only C5). With no hierarchical stage gates in the way, quests finish sooner --- the bot runs through more episodes within the 500 turns --- and that hands the NPCs additional learning opportunities. The implication is that quest structure shapes the \textit{rate} of learning more than its \textit{direction}: both flat and hierarchical variants head toward cooperation, but a flat MDP simply supplies more training data per unit time.

\subsection{Finding 5: Dynamic Lambda Amplifies Cooperation}

C7 (Static $\lambda = 0.325$) achieves a cooperation index of 0.595, compared to C1's 0.699 --- a 0.104 difference, making the dynamic $\lambda$ feedback loop the third-largest contributor to cooperation, after RL itself ($+0.699$) and role masking ($+0.161$). Crucially, C7 still shows strong cooperation emergence (0.000 $\to$ 0.595): the community reward signal alone drives most of the learning, even without the amplifying feedback loop. But the virtuous cycle --- positive community outcomes $\to$ cooperation rises $\to$ $\lambda$ increases $\to$ more weight on community reward --- provides a meaningful additional acceleration. This confirms that the dynamic $\lambda$ mechanism is not merely cosmetic; it is a measurable amplifier of the system's cooperation dynamics.

\begin{figure}[H]
    \centering
    \includegraphics[width=0.85\textwidth]{images/fig_final_bar.png}
    \caption{Final cooperation index by condition after 500 turns. C7 (Static $\lambda$) falls between C6 and C3, confirming the feedback loop's amplifying role.}
    \label{fig:final_bar}
\end{figure}

\section{Reward Decomposition}
\label{sec:reward-decomposition}

\begin{figure}[H]
    \centering
    \includegraphics[width=\textwidth]{images/fig_reward.png}
    \caption{Average reward components over time for C1 (Full System). Community reward C(t) dominates the total signal as $\lambda$ increases. Individual reward G(t) declines as NPC stats stabilize. The famine shock at turn 166 is visible as a dip in community reward.}
    \label{fig:reward}
\end{figure}

The reward decomposition (\cref{fig:reward}) confirms the hypothesized dynamic: early turns are dominated by individual reward $G(t)$ (stat changes from pre-training); as stats stabilize, community reward $C(t)$ becomes the primary gradient signal, amplified by the rising $\lambda$ coefficient.

\section{NPC Adaptation and Role Divergence}
\label{sec:npc-adaptation}

\begin{figure}[H]
    \centering
    \includegraphics[width=\textwidth]{images/fig_npc_adaptation.png}
    \caption{NPC adaptation states after 500 turns. Left: C1 (Full System) shows role-specific divergence. Right: C3 (No RL) learns no cooperation (cooperation tendency stays at 0.000) while the other three coefficients remain at the 0.500 baseline.}
    \label{fig:npc_adaptation}
\end{figure}

\begin{table}[H]
\centering
\caption{NPC Final Adaptation States (C1, Turn 500)}
\label{tab:npc_states}
\begin{tabular}{llccccc}
\toprule
\textbf{NPC} & \textbf{Role} & \textbf{Coop.} & \textbf{Risk} & \textbf{Social} & \textbf{Resil.} & $\boldsymbol{\lambda}$ \\
\midrule
Elder Maren & Elder & 0.697 & 0.488 & 0.899 & 0.500 & 0.433 \\
Farmer Jak & Farmer & 0.699 & 0.413 & 0.763 & 0.500 & 0.434 \\
Aldric & Guard & 0.699 & 0.497 & 0.688 & 0.500 & 0.434 \\
Bryn & Guard & 0.699 & 0.500 & 0.688 & 0.500 & 0.434 \\
Tessa & Tavern Keeper & 0.699 & 0.377 & 0.899 & 0.500 & 0.434 \\
Old Petra & Villager & 0.703 & 0.416 & 0.877 & 0.500 & 0.437 \\
\bottomrule
\end{tabular}
\end{table}

Key observations from \cref{tab:npc_states}:
\begin{itemize}[nosep]
    \item \textbf{Cooperation converges uniformly} --- all NPCs reach $\sim$0.699, indicating the community reward provides a consistent gradient regardless of role.
    \item \textbf{Social sensitivity diverges by role} --- Elder (0.899) and Tavern Keeper (0.899) develop the highest social sensitivity (multiplied by 1.4$\times$ and 1.2$\times$ respectively), with Old Petra close behind (0.877), while the Guards (0.688) remain lowest (0.5$\times$ multiplier) and Farmer Jak sits in between (0.763). This matches role expectations.
    \item \textbf{Risk aversion shows clear variation} --- Tessa (0.377) is least risk-averse, with Farmer Jak (0.413) and Old Petra (0.416) also well below the 0.500 baseline, driven by role-specific risk multipliers encountering occasional negative individual rewards; the Guards stay near baseline (Bryn 0.500, Aldric 0.497).
\end{itemize}

\section{Action Distribution Shift}
\label{sec:action-shift}

\begin{figure}[H]
    \centering
    \includegraphics[width=0.85\textwidth]{images/fig_action_shift.png}
    \caption{Change in NPC action proportions between early and late game (C1). Green = actions used more in late game; red = used less. NPCs shift from social actions (talk, trade) toward utility and exploration (examine, work, rest).}
    \label{fig:action_shift}
\end{figure}

\section{Narrative Quality}
\label{sec:narrative-quality}

\begin{table}[H]
\centering
\caption{Narrative Quality: LLM (qwen3:4b) vs Templates (20 turns)}
\label{tab:narrative}
\begin{tabular}{lccc}
\toprule
\textbf{Metric} & \textbf{C1 (LLM)} & \textbf{C2 (Template)} & \textbf{Winner} \\
\midrule
Contextual grounding & 0.567 & 0.488 & LLM \\
Vocabulary diversity (TTR) & 0.323 & 0.277 & LLM \\
Unique words & 396 & 151 & LLM \\
Avg.\ narration length & 61.4 words & 27.3 words & LLM \\
Lexical coherence & 0.202 & 0.213 & Template \\
\bottomrule
\end{tabular}
\end{table}

The LLM produces narration that is more contextually grounded (0.57 vs 0.49) and more lexically diverse (2.6$\times$ more unique words), while templates score slightly higher on lexical coherence due to reuse of identical phrases.

\section{Social Welfare}
\label{sec:social-welfare}

\begin{figure}[H]
    \centering
    \includegraphics[width=\textwidth]{images/fig_welfare.png}
    \caption{Village social welfare index over 500 turns. Tracks cooperation emergence at the village level, weighted by community reward performance.}
    \label{fig:welfare}
\end{figure}

% =============================================================================
% CHAPTER 6: TOOLS, TECHNOLOGIES, AND PARAMETERS
% =============================================================================
\chapter{Tools, Technologies, and Parameters}
\label{chap:tools}

This chapter documents the technology stack, NPC configuration, and key hyperparameters used in the implementation. The system comprises approximately 17,700 lines of Python (backend) and 6,300 lines of JavaScript (frontend).

\section{Tools and Technologies}
\label{sec:tools-tech}

\begin{itemize}
    \item \textbf{Programming Languages / Frameworks:} Python 3.12 (backend), Vanilla JavaScript (frontend), FastAPI (web framework), Pydantic (data validation).
    \item \textbf{RL Engine:} NumPy for tabular Q-table operations.
    \item \textbf{NLP:} spaCy for input parsing fallback.
    \item \textbf{LLM Integration:} HTTP adapter supporting Ollama, OpenRouter, and LM Studio (OpenAI-compatible endpoints).
    \item \textbf{Frontend Visualization:} Chart.js (analytics dashboard), Cytoscape.js (real-time MDP graph).
    \item \textbf{Research Analytics:} Matplotlib for experiment figures and ablation plots.
    \item \textbf{Testing:} pytest (91 tests across 5 suites: core gameplay, shock engine, adaptation \& masking, analytics, integration).
    \item \textbf{Package Management:} \texttt{uv} with \texttt{pyproject.toml}.
\end{itemize}

\section{NPC Roster and Configuration}
\label{sec:npc-roster}

\begin{table}[H]
\centering
\caption{NPC Roster and Reward Weights}
\label{tab:npc_roster}
\begin{tabular}{llllllll}
\toprule
\textbf{NPC} & \textbf{Role} & \textbf{Location} & \textbf{Happiness} & \textbf{Income} & \textbf{Health} & \textbf{Rep.} \\
\midrule
Elder Maren & Elder & Elder's House & 0.3 & 0.1 & 0.2 & 0.4 \\
Farmer Jak & Farmer & Fields & 0.2 & 0.4 & 0.2 & 0.2 \\
Aldric & Guard & Gate & 0.1 & 0.2 & 0.3 & 0.4 \\
Bryn & Guard & Gate & 0.1 & 0.2 & 0.3 & 0.4 \\
Tessa & Tavern Keeper & Tavern & 0.2 & 0.4 & 0.1 & 0.3 \\
Old Petra & Villager & Village Center & 0.4 & 0.2 & 0.2 & 0.2 \\
\bottomrule
\end{tabular}
\end{table}

\section{Key Parameters and Settings}
\label{sec:parameters}

\begin{table}[H]
\centering
\caption{Important Parameters Used in Implementation}
\label{tab:parameters}
\begin{tabularx}{\textwidth}{|l|X|}
\hline
\textbf{Parameter} & \textbf{Value / Description} \\
\hline
Learning rate ($\alpha$) & 0.1 --- standard tabular Q-learning rate \\
\hline
Discount factor ($\gamma$) & 0.9 --- balances immediate vs future reward \\
\hline
Exploration ($\varepsilon$) & Decays from 0.15 to 0.05 at rate 0.995/turn \\
\hline
Pre-training & 100 episodes $\times$ 50 turns per NPC at $\varepsilon = 0.5$ \\
\hline
Adaptation rate & 0.018 base rate with homeostatic drift 0.004 toward 0.5 \\
\hline
Dynamic $\lambda$ & $\lambda_i = 0.05 + \text{coop}_i \times 0.55$, range [0.05, 0.60] \\
\hline
Static $\lambda$ (C7) & Fixed at 0.325 (midpoint of dynamic range) \\
\hline
Soft mask bonus/penalty & $+0.5$ for role-aligned actions, $-0.3$ for misaligned \\
\hline
HP penalty threshold & $-10$ reward when HP drops below 20\% \\
\hline
Nudge reward & $R_{\text{nudge}} = 0.3 \cdot (d_{\text{before}} - d_{\text{after}})$ \\
\hline
Master seed & 42 (all randomness derived deterministically) \\
\hline
Experiment duration & 500 turns per condition with quest auto-restart \\
\hline
Famine shock injection & Turn 166 (one-third of session) \\
\hline
\end{tabularx}
\end{table}

% =============================================================================
% CHAPTER 7: CONCLUSION AND FUTURE SCOPE
% =============================================================================
\chapter{Conclusion and Future Scope}
\label{chap:conclusion}

This project demonstrates that emergent cooperative behavior can arise from independent Q-learning agents when the reward signal includes a community welfare component gated by a dynamic $\lambda$ coefficient. The key findings:

\begin{enumerate}
    \item \textbf{Cooperation emerges from learning.} $\mathcal{K}(t)$ rises from 0.000 to 0.699 over 500 turns with Q-learning, while schedule-only NPCs show no change ($+0.699$ gap). This confirms the research hypothesis and makes RL by far the dominant driver.
    \item \textbf{Role masking provides the strongest inductive bias.} A 0.161 improvement in cooperation when NPCs have role-constrained action spaces, producing cleaner learning signals --- the largest single-component contribution after RL itself.
    \item \textbf{Shocks impose a lasting cost, yet agents keep learning.} Removing shocks raises final cooperation by 0.237 (C5 reaches 0.936 vs.\ C1's 0.699); the famine at turn 166 lastingly depresses C1's trajectory. The agents do not collapse --- cooperation keeps rising afterward --- but they do not recover the lost ground within the session.
    \item \textbf{The dynamic $\lambda$ feedback loop is a measurable amplifier.} Static $\lambda$ still produces strong emergence (0.000 $\to$ 0.595), but dynamic $\lambda$ adds $+0.104$ through the virtuous cycle --- the third-largest single-component contribution, after RL and role masking.
    \item \textbf{LLM narration adds quality without losing context.} With tuned prompts, the LLM achieves higher contextual grounding than templates while producing 2.6$\times$ richer vocabulary.
\end{enumerate}

\section{Limitations}
\label{sec:limitations}

\begin{itemize}[nosep]
    \item Results are from automated bot playthroughs with a single strategy (\texttt{quest\_focused}); human player behavior may produce different cooperation dynamics.
    \item The adaptation rate (0.018) and role-specific multipliers were empirically tuned; systematic sensitivity analysis is needed.
    \item Tabular Q-learning limits scalability --- 225 states is manageable for 6 NPCs but would not scale to larger worlds without function approximation.
    \item Narrative coherence metrics rely on token overlap rather than semantic similarity; sentence-transformer embeddings would provide stronger signal.
\end{itemize}

\section{Future Scope}
\label{sec:future-scope}

\begin{itemize}[nosep]
    \item Scale to 100 bot playthroughs per condition with varied strategies for statistical significance.
    \item Conduct user studies (20 participants per condition) with post-session questionnaires.
    \item Compare tabular Q-learning vs.\ DQN for NPC behavior.
    \item Extend to branching quest graphs with multiple concurrent quest lines.
    \item Explore longer sessions ($>$1000 turns) to test cooperation saturation.
\end{itemize}

\section{Real-World Transfer}
\label{sec:real-world-transfer}

The decomposed reward framework $R_i(t) = P_i(t) + G_i(t) + \lambda_i(t) \cdot C(t)$ generalises beyond games to any domain where heterogeneous subsystems must balance local objectives against a shared global target.

\textbf{Smart Factory Coordination.}
In an automated production line, each cell (stamping, welding, painting) optimises a local KPI---throughput, defect rate, idle time---analogous to role-specific $G_i$ weights. The community signal $C(t)$ maps to plant-wide OEE (Overall Equipment Effectiveness). Supply-chain disruptions, equipment failures, and demand spikes correspond directly to famine, plague, and trade-boom shocks. Current Industry~4.0 approaches rely on centralised MPC controllers that must be re-modelled when the line changes; a dynamic-$\lambda$ architecture lets each cell run a lightweight learner and join the reward structure without replanning the plant. Our C1~vs.~C3 result ($+0.699$ cooperation from RL over fixed schedules) suggests learned policies would significantly outperform static dispatch rules.

\textbf{Automotive ECU Self-Tuning.}
A modern engine's sensor--actuator subsystems (knock sensor, O$_2$ sensor, throttle body, MAP, coolant temperature, crank position) map onto heterogeneous NPC roles: each maximises a local objective (detonation margin, fuel-trim accuracy, airflow response) while the community target $C(t)$ is engine-level efficiency (e.g.\ brake-specific fuel consumption). The dynamic $\lambda$ feedback loop enables continuous self-calibration: a factory base map provides the pre-trained Q-table, and on-road experience adjusts cooperation tendencies as altitude, fuel quality, and component wear shift the reward landscape. Component aging becomes a slow-onset shock; aftermarket modifications (larger turbo, free-flowing exhaust) reset part of the state, and the system re-converges without manual tuning---our C4 result ($\mathcal{K} = 0.871$) suggests rapid re-convergence is achievable.

\textbf{Bridging the Gap.}
Two challenges must be addressed: (i)~the 225-state tabular formulation must give way to function approximation (DQN or policy gradients) for high-dimensional, real-time sensor streams; and (ii)~the $P_i$ penalty discourages but does not formally prevent catastrophic states---a hybrid architecture enforcing hard constraints from classical control theory while allowing RL-based optimisation within the safe envelope is the likely path forward.

% =============================================================================
% BIBLIOGRAPHY
% =============================================================================
\addcontentsline{toc}{chapter}{Bibliography}
\bibliographystyle{plainnat}
\bibliography{references}

\end{document}